% chapter01.tex - 第一章：绪论
% 使用了 \label{} 命令创建标签，便于交叉引用
% 使用了 itemize 和 enumerate 列表环境

\chapter{chp1 绪论}

\begin{itemize}
    \item \textbf{医疗诊断}：辅助医生进行疾病筛查和诊断
    \item \textbf{工业检测}：产品质量自动检测
\end{itemize}

\begin{enumerate}
    \item AlexNet（2012年）—— 开创了深度学习在图像识别中的应用
    \item Transformer（2020年）—— 基于注意力机制的视觉模型
\end{enumerate}

\section{整体框架}
本文提出的方法整体框架如图\ref{fig:framework}所示。

\begin{figure}[htbp]
    \centering
    % \includegraphics[width=0.8\textwidth]{figures/framework.png}
    \caption{方法整体框架图}
    \label{fig:framework}
\end{figure}

\section{算法描述}
本文提出的算法如算法 \ref{alg:main} 所示。

\begin{algorithm}[htbp]
    \caption{主要算法}
    \label{alg:main}
    \begin{algorithmic}[1]
        \REQUIRE 输入数据 $X$
        \ENSURE 输出结果 $Y$
        \STATE 初始化参数
        \FOR{$i = 1$ to $N$}
            \STATE 处理步骤1
            \STATE 处理步骤2
        \ENDFOR
        \RETURN $Y$
    \end{algorithmic}
\end{algorithm}

这里是引用1\cite{krizhevsky2012imagenet}

这里是引用2\cite{vaswani2017attention}  

\begin{equation}
    y = f\left(\sum_{i=1}^{n} w_i x_i + b\right)
\end{equation}
其中，$x_i$为输入，$w_i$为权重，$b$为偏置，$f(\cdot)$为激活函数。

\begin{table}[htbp]
    \centering
    \caption{常用图像分类算法比较}
    \label{tab:classification-algorithms}
    \begin{tabular}{p{0.25\textwidth} p{0.35\textwidth} p{0.3\textwidth}}
        \toprule
        \textbf{算法类型} & \textbf{优点} & \textbf{缺点} \\
        \midrule
        支持向量机（SVM） & 小样本效果好，泛化能力强 & 大规模数据训练慢 \\
        随机森林（RF） & 抗过拟合，可处理高维数据 & 模型解释性差 \\
        卷积神经网络（CNN） & 自动特征提取，识别精度高 & 需要大量训练数据 \\
        \bottomrule
    \end{tabular}
\end{table}

\subsection{精确率（Precision）和召回率（Recall）}
\begin{align}
    \text{Precision} &= \frac{TP}{TP + FP} \label{eq:precision} \\
    \text{Recall} &= \frac{TP}{TP + FN} \label{eq:recall}
\end{align}